\documentclass[10pt,a4paper]{amsart}
\usepackage[margin=19mm]{geometry}
\usepackage{amsmath,amssymb,mathtools}
\usepackage[scr=rsfs,cal=euler]{mathalfa}
\usepackage{xfrac}
\usepackage[unicode,hidelinks,bookmarks=false]{hyperref}
\newcommand{\ie}{i.e.\ }
\newcommand{\cf}{cf.\ }
\newcommand{\resp}{respectively\ }
\newcommand{\Subsection}{\subsection}
\providecommand{\nobreakdash}{\nobreak\hbox{-}}
\setlength{\emergencystretch}{2em}
\multlinegap0pt
\usepackage{tikz}
\usepackage{xcolor}
\usepackage{amssymb}
\usepackage{mathtools}
\usepackage[all]{xy}
\usepackage{cancel}
\usepackage[normalem]{ulem}
\newtheorem{prop}{Proposition}
\newtheorem{claim}{Claim}
\DeclareMathOperator{\Base}{\mathsf{Base}}
\newcommand{\cA}{\mathcal{A}}
\newcommand{\cB}{\mathcal{B}}
\newcommand{\cC}{\mathcal{C}}
\newcommand{\cM}{\mathcal{M}}
\newcommand{\cO}{\mathcal{O}}
\newcommand{\cV}{\mathcal{V}}
\newcommand{\cW}{\mathcal{W}}
\title{Open-closed maps and spectral local systems}
\author{Noah Porcelli, Ivan Smith}
\date{Source: arXiv:2509.21483v2 (CC BY 4.0)}
\begin{document}
\maketitle

\noindent\emph{Source and licence.} Open-closed maps and spectral local systems. Version arXiv:2509.21483v2 (CC BY 4.0), distributed by arXiv under the Creative Commons Attribution 4.0 International licence, http://creativecommons.org/licenses/by/4.0/, as recorded in the arXiv licence metadata for this exact version and shown on the abstract page https://arxiv.org/abs/2509.21483v2. Full licence notice: \texttt{licenses/arxiv-2509.21483-CC-BY-4.0.txt}.\par\smallskip

\noindent\textit{Editorial scope.} Counted unit: the article's prop modification (source0.tex lines 2717--2721). The article's complete original proof of it is delivered with it (source0.tex lines 2731--2773, one \texttt{proof} environment of 43 lines, which the article places away from the statement). No mathematics is written, repaired or re-derived: the statement and the proof are the article's own lines, in the article's own notation, and no retained line is substituted or re-flowed.  Retained uncounted as the source-local context the proof needs: lines 2683--2691.  Nothing was omitted from any retained span, so the package carries no omission record.  No retained line contains a citation, so the wrapper carries no bibliography and no bibliography entry.  No retained line contains a figure environment, an image-inclusion command or drawing code, so no image is delivered and the package declares no asset dependency.  Editorial formatting only, in addition: an AMS \texttt{amsart} preamble that restates the article's own declarations and macros used by the retained lines, namely the environments \texttt{prop}, \texttt{claim} on separate counters, together with the standard packages those lines need.  The retained environments print in this document as Proposition 1, Claim 1 in order of appearance, whereas the article prints the same items with the numbering of its own full document.  The complete native source of the article as distributed in the arXiv e-print for this exact version is bundled unchanged as \texttt{sources/185747/source0.tex}.
    \begin{prop}\label{prop: lift ob}
        Let $\cA$ and $\cB$ be $E$-oriented $\tau_{\leq i+1}$-flow categories, and $\cW: \cA \to \cB$ an $E$-oriented $\tau_{\leq i}$-morphism. Then there is a homology class

        \begin{equation}\label{eq: lift ob}
            [\cO^E] \in \operatorname{Hom}\left(CM_{*+i+2}(\cA), CM_*(\cB; \Omega^{E}_{i+1})\right)
        \end{equation}
        such that if $[\cO^E]$ vanishes, $\cW$ is the truncation of an $E$-oriented $\tau_{\leq i+1}$-morphism $\cW'$.
    
    \end{prop}

    \begin{prop}\label{prop: truncation modification}
        Let $\cA, \cB$ be $E$-oriented $\tau_{\leq i+1}$-flow categories, and $\cW: \cA \to \cB$ a $E$-oriented $\tau_{\leq i}$-morphism.

        Then there is a $E$-oriented $\tau_{\leq i+1}$-flow category $\cC$, and a $E$-oriented $\tau_{\leq i+1}$-equivalence $\cV: \cA \to \cC$, and such that $\tau_{\leq i} \cC = \tau_{\leq i} \cB$ and $\tau_{\leq i} \cV = \cW$.
    \end{prop}

    \begin{proof}[Proof of Proposition \ref{prop: truncation modification}]
        $\cC$ is defined to have objects the same as $\cB$ (with the same gradings), and
        \begin{equation}
            \cC_{yy'} = \begin{cases}
                \cB_{yy'} 
                &
                \textrm{ if } |y|-|y'|-1 \leq i
                \\
                \cB_{yy'} \sqcup P_{yy'}
                &
                \textrm{ if } |y|-|y'|-1 = i+1
            \end{cases}
        \end{equation}
        where for each $y, y' \in \cC$ with $|y|-|y'|-1 = i+1$, $P_{yy'}$ is a choice of closed $(i+1)$-manifold with tangent bundle stably classified by a map $P_{yy'} \to \Base(E)$ (we will choose the $P_{yy'}$ in due course). $\cC$ is a $E$-oriented pre-$\tau_{\leq i+1}$ flow category, where the maps to $E$ are inherited from those given on $\cB$ and the $P_{yy'}$. We write this as $\cC=\cC(P)$ to indicate the choice involved; in particular, $\cB = \cC(\emptyset)$.

        By taking a matrix whose coefficients are given by the bordism class of each $P_{yy'}$, we obtain a linear map of chain complexes
        \begin{equation}
            P: CM_{*+i+2}(\cB) \to CM_*(\cB; \Omega_{i+1}^{E})
        \end{equation}
        

        \begin{claim}
            $\cC$ is an $E$-oriented $\tau_{\leq i+1}$-flow category if and only if $P$ is a map of chain complexes.
        \end{claim}
        \begin{proof}[Proof of claim.]

            Let $(dP)_{xy} \in \Omega^{E}_{i+1}$ be the coefficients of the matrix representing the linear map
            \begin{equation}
                dP: CM_{*+i+3}(\cB) \to CM_*\left(\cB; \Omega_{i+1}^{E}\right)
            \end{equation}
            By construction, a nullbordism of (representatives of) all $(dP)_{xy}$ is exactly the data of an extension of the pre-$\tau_{\leq i+1}$-flow category $\cC$.
        \end{proof}

        $\cW$ defines a $E$-oriented $\tau_{\leq i}$-morphism $\cA \to \cC$. From Proposition \ref{prop: lift ob}, the obstruction to this admitting an extension to a $E$-oriented $\tau_{\leq i+1}$-morphism $\cV: \cA \to \cC$ is a class
        \begin{equation}
            [\cO^E(P)] \in \operatorname{Hom}\left(CM_{*+i+2}(\cA), \cM_{*}(\cB; \Omega^{E}_{i+1}) \right)
        \end{equation}
        Inspecting the construction of the class $[\cO^E(P)]$, we find that
        \begin{equation}
            [\cO^E(P)] = [\cO^E(0) + P \circ \cW_*]
        \end{equation}
        where $\cW_*: CM_*(\cA) \to CM_*(\cB)$ is the map induced by $\cW$. Since $\cW_*$ is a quasi-isomorphism, we may therefore choose $P$ so that the homology class of $[\cO^E(P)]$ vanishes, and so the $\Psi$-oriented $\tau_{\leq i+1}$-extension $\cV$ of $\cW$ exists for this $\cC=\cC(P)$.
    \end{proof}

\end{document}
